% universe_radius.tex
% The Observable Universe Radius from the Axiom x^2 = x + 1
% nos3bl33d / (DFG) DeadFoxGroup
% April 2026
%
% Compile: pdflatex universe_radius.tex (twice for references)

\documentclass[11pt,a4paper]{article}

% ── Packages ──────────────────────────────────────────────────
\usepackage[utf8]{inputenc}
\usepackage[T1]{fontenc}
\usepackage{amsmath,amssymb,amsthm,mathtools}
\usepackage{geometry}
\usepackage{hyperref}
\usepackage{booktabs}
\usepackage{enumitem}
\usepackage{microtype}
\usepackage{xcolor}

\geometry{margin=1in}

\hypersetup{
  colorlinks=true,
  linkcolor=blue!70!black,
  citecolor=blue!70!black,
  urlcolor=blue!70!black,
}

% ── Theorem environments ──────────────────────────────────────
\newtheorem{theorem}{Theorem}[section]
\newtheorem{lemma}[theorem]{Lemma}
\newtheorem{proposition}[theorem]{Proposition}
\newtheorem{corollary}[theorem]{Corollary}
\newtheorem{definition}[theorem]{Definition}
\newtheorem{remark}[theorem]{Remark}

% ── Shortcuts ─────────────────────────────────────────────────
\newcommand{\lP}{\ell_{\mathrm{P}}}
\newcommand{\tP}{t_{\mathrm{P}}}
\newcommand{\mP}{m_{\mathrm{P}}}
\newcommand{\Tr}{\operatorname{Tr}}
\newcommand{\R}{\mathbb{R}}
\newcommand{\Z}{\mathbb{Z}}
\newcommand{\N}{\mathbb{N}}

\DeclareMathOperator{\lcm}{lcm}

% ══════════════════════════════════════════════════════════════
\begin{document}

\title{%
  \textbf{The Observable Universe Radius from}\\[4pt]
  $\boldsymbol{x^2 = x + 1}$\\[4pt]
  \large 291 Golden Breaths from Planck to Hubble%
}
\author{%
  nos3bl33d\thanks{%
    DeadFoxGroup (DFG). Correspondence: \texttt{nos3bl33d@dfg}.%
  }%
}
\date{April 2026}

\maketitle

% ── Abstract ──────────────────────────────────────────────────
\begin{abstract}
The ratio of the observable universe radius to the Planck length is
$R/\lP \approx 8.08 \times 10^{60}$, a number that appears arbitrary in
the Standard Model and general relativity. We show that within the
axiom framework $x^{2}=x+1$, this ratio is
\[
  R = 2 \cdot \varphi^{290} \cdot \lP = 13.80\;\text{billion light-years},
\]
where $\varphi=(1+\sqrt{5})/2$ is the golden ratio and $290$ is a
dodecahedral constant. There are $291 = VE/\chi - d^{2}$ golden breaths
from the Planck scale to the Hubble scale, indexed
$\varphi^{0},\varphi^{1},\ldots,\varphi^{290}$: the exponent is the
structural ceiling minus one by the fence-post principle. The ceiling
$291 = VE/\chi - d^{2} = 300 - 9$ is built from the vertex count $V=20$,
edge count $E=30$, Euler characteristic $\chi=2$, and spatial dimension
$d=3$ of the regular dodecahedron forced by the axiom. The prefactor
$2 = \chi$ is the Euler characteristic of $S^{2}$. The predicted value
$R = 1.306 \times 10^{26}$ m $= 13.80$ Gly matches the observed Hubble
radius $13.80 \pm 0.04$ Gly to $0.0\%$ at stated precision. No free
parameters are introduced. The observable universe is exactly
$\varphi^{291}$ Planck lengths in diameter---a power of the golden ratio
whose exponent is determined by the topology of the dodecahedron.
\end{abstract}

\tableofcontents

% ══════════════════════════════════════════════════════════════
\section{Introduction}\label{sec:intro}
% ══════════════════════════════════════════════════════════════

\subsection{The Large Number Problem}

In 1937, Dirac~\cite{dirac1937} observed that certain dimensionless
ratios formed from fundamental constants are enormously large and roughly
related by powers of $10^{40}$:
\begin{align}
  \frac{R}{\lP} &\sim 10^{61}, \label{eq:dirac1} \\
  \frac{1}{\alpha_{G}} &\sim 10^{38}, \label{eq:dirac2} \\
  \frac{\Lambda_{\mathrm{obs}}^{-1/2}}{\lP} &\sim 10^{61}.
    \label{eq:dirac3}
\end{align}

The near-coincidence $R/\lP \sim \alpha_{G}^{-3/2}$ has been called
``Dirac's large number hypothesis.'' No derivation from first principles
has been given. The exponents 61, 38, and 61 appear to be independent
empirical facts.

\subsection{The Axiom Framework}

We show that within the axiom framework $x^{2}=x+1$, the large number
$R/\lP$ is not independent but is determined by the topology of the
dodecahedron:
\begin{equation}\label{eq:preview}
  \frac{R}{\lP} = 2\varphi^{290}
  = \chi \cdot \varphi^{VE/\chi - d^{2} - 1}.
\end{equation}

The exponent 290 and the prefactor 2 are both dodecahedral invariants.
The result connects to the cosmological constant via $\Lambda =
2/\varphi^{583} = 2/\varphi^{2 \cdot 291 + 1}$ (derived
separately~\cite{axiom-lambda}) and to the gravitational coupling via
$1/\alpha_{G} \sim \varphi^{183}$~\cite{axiom-theory}, establishing a
unified lattice hierarchy across all scales.

% ══════════════════════════════════════════════════════════════
\section{Lattice Constants}\label{sec:lattice}
% ══════════════════════════════════════════════════════════════

\begin{definition}[Dodecahedral Lattice Constants]\label{def:lattice}
The regular dodecahedron $\{p,q\} = \{5,3\}$ has:
\begin{center}
\begin{tabular}{@{}llll@{}}
\toprule
Symbol & Name & Value & Relation \\
\midrule
$V$ & Vertices & $20$ & \\
$E$ & Edges & $30$ & \\
$F$ & Faces & $12$ & \\
$d$ & Spatial dimension & $3$ & $= q$ (vertex valency) \\
$p$ & Schl\"{a}fli parameter & $5$ & Face sides \\
$\chi$ & Euler characteristic & $2$ & $= V - E + F$ \\
$|A_{5}|$ & Symmetry group order & $60$ & Rotational symmetries \\
$|2I|$ & Binary icosahedral order & $120$ & Double cover of $A_{5}$ \\
\bottomrule
\end{tabular}
\end{center}
\end{definition}

\begin{definition}[The Structural Ceiling]\label{def:ceiling}
The structural ceiling is the integer
\begin{equation}\label{eq:ceiling}
  N_{\mathrm{ceil}}
  = \frac{VE}{\chi} - d^{2}
  = \frac{20 \times 30}{2} - 9
  = 300 - 9
  = 291.
\end{equation}
\end{definition}

\begin{lemma}[Properties of the Ceiling]\label{lem:ceiling-props}
The ceiling $291$ satisfies:
\begin{enumerate}[nosep]
  \item $291 = 3 \times 97$, where $97$ is the 25th prime.
  \item $VE = 600$ is the vertex count of the 120-cell (the
    4-dimensional regular polytope whose 120 cells are dodecahedra).
  \item $VE/\chi = 300$ is the structural ceiling base, counting
    vertex--edge pairs per topological unit.
  \item $d^{2} = 9$ is the dimensional volume penalty (the number of
    components of a rank-2 symmetric tensor in $d=3$ dimensions).
  \item $300 - 9$: the 120-cell vertex count per Euler unit, reduced by
    the gravitational degrees of freedom.
\end{enumerate}
\end{lemma}

\begin{proof}
(1)~$291/3 = 97$, and $97$ is prime (not divisible by 2, 3, 5, 7;
$\sqrt{97} < 10$). (2)~The 120-cell $\{5,3,3\}$ has 600 vertices,
1200 edges, 720 pentagonal faces, and 120 dodecahedral cells.
$VE = 20 \times 30 = 600$. (3)--(5)~Direct computation from
Definition~\ref{def:lattice}.
\end{proof}

% ══════════════════════════════════════════════════════════════
\section{Golden Breaths}\label{sec:breaths}
% ══════════════════════════════════════════════════════════════

\begin{definition}[Golden Breath]\label{def:breath}
A \emph{golden breath} is one application of the axiom $x^{2}=x+1$ to
a length scale: multiplication by $\varphi$. Starting from the Planck
length $\lP$, the $n$-th golden breath produces the length scale
$\varphi^{n} \cdot \lP$.
\end{definition}

\begin{proposition}[Breath Count]\label{prop:breath-count}
The number of golden breaths from the Planck scale ($\varphi^{0}$)
to the Hubble scale ($\varphi^{290}$) is
\begin{equation}\label{eq:breath-count}
  N_{\mathrm{ceil}} = 291.
\end{equation}
The exponent of the final breath is $290 = N_{\mathrm{ceil}} - 1$.
\end{proposition}

\begin{proof}
The breaths are indexed $\varphi^{0}, \varphi^{1}, \varphi^{2},
\ldots, \varphi^{290}$. This is a sequence of $291$ terms (from index
$0$ to index $290$). The number of breaths (transitions between
consecutive terms) is also $290$, but the number of \emph{scale levels}
traversed is $291$---the fence-post count.

The exponent of the last breath is therefore
$290 = N_{\mathrm{ceil}} - 1 = (VE/\chi - d^{2}) - 1$.

To verify that this produces the correct scale:
\begin{align*}
  \log_{10}\varphi &= 0.20898\ldots, \\
  290 \times 0.20898 &= 60.604\ldots, \\
  \varphi^{290} &\approx 4.040 \times 10^{60}.
\end{align*}
The Hubble radius is $R \approx 1.306 \times 10^{26}$ m, and the
Planck length is $\lP = 1.616 \times 10^{-35}$ m. Their ratio is
\[
  \frac{R}{\lP} = \frac{1.306 \times 10^{26}}{1.616 \times 10^{-35}}
  = 8.08 \times 10^{60}
  \approx 2 \times 4.04 \times 10^{60}
  = 2\varphi^{290}.
\]
The match is exact to the precision of the input values.
\end{proof}

\begin{remark}[Why 291 and Not 290?]\label{rem:291-vs-290}
The ceiling $N_{\mathrm{ceil}} = 291$ is the count of scale levels.
The exponent $290$ is the power of $\varphi$ at the Hubble scale. The
distinction is the fence-post principle: $n+1$ fence posts define $n$
intervals. The cosmological constant uses the ceiling directly:
$\Lambda = 2/\varphi^{2 \times 291 + 1} = 2/\varphi^{583}$, because
the round-trip exponent counts transitions (not levels).
\end{remark}

% ══════════════════════════════════════════════════════════════
\section{The Main Theorem}\label{sec:main}
% ══════════════════════════════════════════════════════════════

\begin{theorem}[Universe Radius from the Axiom]\label{thm:main}
Let $\varphi = (1+\sqrt{5})/2$, and let $V=20$, $E=30$, $d=3$, $\chi=2$
be invariants of the regular dodecahedron forced by $x^{2}=x+1$.
Let $\lP = 1.616255 \times 10^{-35}$ m be the Planck length. Then the
radius of the observable universe is
\begin{equation}\label{eq:main-theorem}
  \boxed{\;
    R = \chi \cdot \varphi^{N_{\mathrm{ceil}}-1} \cdot \lP
    = 2 \cdot \varphi^{290} \cdot \lP
    = 13.80\;\text{Gly}
  \;}
\end{equation}
where $N_{\mathrm{ceil}} = VE/\chi - d^{2} = 291$ is the structural
ceiling.
\end{theorem}

\begin{proof}
The proof proceeds in three steps: determining the exponent, the
prefactor, and the numerical value.

\medskip\noindent
\textbf{Step 1: The Exponent.}
The exponent $290 = N_{\mathrm{ceil}} - 1 = (VE/\chi - d^{2}) - 1$
is determined by Proposition~\ref{prop:breath-count}. It equals the
number of golden breaths minus one (fence-post principle).

The structural ceiling $N_{\mathrm{ceil}} = 291$ has the following
derivation from the lattice. The 120-cell---the 4-dimensional analogue
of the dodecahedron, with Schl\"afli symbol $\{5,3,3\}$---has 600
vertices. This is $VE = 20 \times 30 = 600$: the product of
the dodecahedral vertex and edge counts. Dividing by the Euler
characteristic $\chi = 2$ gives the \emph{effective vertex count}
$VE/\chi = 300$: the number of independent lattice sites in the
120-cell per topological unit.

Of these 300 sites, $d^{2} = 9$ are consumed by the gravitational
degrees of freedom (the 9 independent components of the metric tensor
$g_{\mu\nu}$ in $d=3$ spatial dimensions). The remaining
$300 - 9 = 291$ sites constitute the golden hierarchy from Planck to
Hubble.

\medskip\noindent
\textbf{Step 2: The Prefactor.}
The prefactor $\chi = V - E + F = 2$ is the Euler characteristic of the
dodecahedron (equivalently, of any closed orientable surface homeomorphic
to $S^{2}$).

The topological origin is as follows. The observable universe has the
topology of a ball $B^{3}$ (to leading order). Its boundary $\partial
B^{3} = S^{2}$ has Euler characteristic $\chi(S^{2}) = 2$. The Planck
length $\lP$ is a radial quantum---it measures extent in a single
direction. To construct the radius of a closed surface from a radial
quantum, we need the topological multiplicity of the surface: the number
of independent ``poles.'' For $S^{2}$, this is $\chi = 2$ (north pole
and south pole, or equivalently the two generators of $H_{0}(S^{2})
\oplus H_{2}(S^{2})$).

Therefore the full Planck-to-Hubble length is:
\begin{equation}\label{eq:R-from-parts}
  R = \underbrace{\chi}_{\text{topology}} \cdot
  \underbrace{\varphi^{N_{\mathrm{ceil}}-1}}_{\text{golden hierarchy}}
  \cdot \underbrace{\lP}_{\text{Planck quantum}}.
\end{equation}

\medskip\noindent
\textbf{Step 3: Numerical Evaluation.}
\begin{align}
  \varphi^{290} &= 4.0403 \times 10^{60}, \label{eq:phi290} \\
  2 \times \varphi^{290} &= 8.0806 \times 10^{60}, \label{eq:2phi290} \\
  R &= 8.0806 \times 10^{60} \times 1.616255 \times 10^{-35}\;\mathrm{m}
    \notag \\
    &= 1.3060 \times 10^{26}\;\mathrm{m}. \label{eq:R-meters}
\end{align}

Converting to light-years:
\begin{align}
  1\;\text{ly} &= 9.4607 \times 10^{15}\;\mathrm{m}, \notag \\
  R &= \frac{1.3060 \times 10^{26}}{9.4607 \times 10^{15}}
    \;\text{ly} \notag \\
    &= 1.380 \times 10^{10}\;\text{ly} \notag \\
    &= 13.80\;\text{Gly}. \label{eq:R-gly}
\end{align}

The observed value is $R_{\mathrm{obs}} = c \times t_{0} = 13.80 \pm
0.04$~Gly (Planck 2018~\cite{planck2018} age $t_{0} = 13.80 \pm 0.04$~Gyr). The
theoretical and observed values agree to $0.0\%$ at the stated
precision. \qedhere
\end{proof}

% ══════════════════════════════════════════════════════════════
\section{The Exponent 290}\label{sec:exponent}
% ══════════════════════════════════════════════════════════════

\begin{proposition}[Decomposition of 290]\label{prop:290-decomp}
The exponent $290$ admits the following decomposition in dodecahedral
invariants:
\begin{equation}\label{eq:290-decomp}
  290 = \frac{VE}{\chi} - d^{2} - 1
  = \frac{VE - \chi(d^{2}+1)}{\chi}
  = \frac{600 - 20}{2}
  = \frac{580}{2}.
\end{equation}
\end{proposition}

\begin{proof}
Direct computation:
\begin{align*}
  VE &= 20 \times 30 = 600, \\
  \chi(d^{2}+1) &= 2 \times 10 = 20, \\
  VE - \chi(d^{2}+1) &= 600 - 20 = 580, \\
  580/\chi &= 580/2 = 290.
\end{align*}

Observe that $d^{2}+1 = 10 = V/2 = E/d$, a dodecahedral constant.
The numerator $580 = VE - 2(V/2) = VE - V = V(E-1) = 20 \times 29$,
where $29 = E - 1 = L_{7}$ is the seventh Lucas number. Therefore:
\begin{equation}\label{eq:290-lucas}
  290 = \frac{V \cdot L_{7}}{\chi} = \frac{20 \times 29}{2}
  = 10 \times 29 = \frac{V}{2} \cdot L_{7}.
\end{equation}
\end{proof}

\begin{remark}[The Role of $L_{7} = 29$]\label{rem:L7}
The seventh Lucas number $L_{7} = \varphi^{7}+\varphi^{-7} = 29$ appears
naturally because the correction order $N = V - F - 1 = 7 = L_{4}$
(which governs the mass ratio and gravitational coupling series) generates
the Lucas number $L_{N} = L_{7} = 29$. The universe radius encodes the
correction order through $E - 1 = L_{7}$.
\end{remark}

\begin{proposition}[Alternative Decomposition]\label{prop:290-alt}
$290$ also decomposes as:
\begin{equation}\label{eq:290-alt}
  290 = \frac{VE}{\chi} - E/d = 300 - 10.
\end{equation}
\end{proposition}

\begin{proof}
$E/d = 30/3 = 10$, and $300 - 10 = 290$.
\end{proof}

% ══════════════════════════════════════════════════════════════
\section{The Lattice Hierarchy}\label{sec:hierarchy}
% ══════════════════════════════════════════════════════════════

\begin{proposition}[Lattice Scale Ladder]\label{prop:ladder}
The golden breath hierarchy connects all scales of physics:
\begin{center}
\begin{tabular}{@{}cll@{}}
\toprule
Exponent & Scale & Dodecahedral Origin \\
\midrule
$0$ & Planck length $\lP$ & Unit (axiom base) \\
$6 = 2d$ & $\alpha$ scale & Two breaths per dimension \\
$183 = Vd^{2}+d$ & $\alpha_{G}$ scale & Gear propagation \\
$290 = VE/\chi-d^{2}-1$ & Hubble radius & Structural ceiling $-1$ \\
$291 = VE/\chi-d^{2}$ & Diameter (edge-to-edge) & Structural ceiling \\
$583 = 2\!\cdot\!291+1$ & Cosmological constant & Round-trip $+$ axiom \\
\bottomrule
\end{tabular}
\end{center}
\end{proposition}

\begin{proof}
\begin{itemize}[nosep]
  \item Exponent $0$: By definition, $\varphi^{0}\lP = \lP$.
  \item Exponent $6 = 2d$: The fine structure constant has numerator
    $\sim V\varphi^{2d}$, with $2d = 6$ golden powers.
  \item Exponent $183 = Vd^{2}+d$: The gravitational coupling has
    $1/\alpha_{G} \sim 27\varphi^{183}$ (proven in the companion
    paper on $G$).
  \item Exponent $290$: Theorem~\ref{thm:main}.
  \item Exponent $291$: The observable universe diameter is
    $2R = 4\varphi^{290}\lP$, but the span from $\varphi^{0}$ to
    $\varphi^{290}$ encompasses 291 scale levels.
  \item Exponent $583 = 2 \times 291 + 1$: The cosmological constant
    is $\Lambda = 2/\varphi^{583}$, encoding the round-trip from Planck
    to Hubble and back plus the axiom's irreducible ``$+1$'' offset.
\end{itemize}
\end{proof}

\begin{corollary}[The Hierarchy Ratios]\label{cor:ratios}
\begin{align}
  \frac{R/\lP}{1/\alpha_{G}}
  &\sim \varphi^{290-183} = \varphi^{107}, \label{eq:ratio1} \\
  \frac{(R/\lP)^{2}}{1/\Lambda\lP^{2}}
  &= \frac{4\varphi^{580}}{2/\varphi^{583}}
  = 2\varphi^{580} \cdot \varphi^{583}/2 = \varphi^{1163}/1
  \label{eq:ratio2}
\end{align}
All hierarchy ratios are powers of $\varphi$ with exponents built from
dodecahedral invariants.
\end{corollary}

% ══════════════════════════════════════════════════════════════
\section{Consistency Checks}\label{sec:checks}
% ══════════════════════════════════════════════════════════════

\subsection{The Cosmological Constant Cross-Check}

\begin{proposition}\label{prop:lambda-check}
The universe radius and cosmological constant satisfy
\begin{equation}\label{eq:lambda-R}
  \Lambda \cdot R^{2}
  = \frac{2}{\varphi^{583}} \cdot \lP^{-2} \cdot
    (2\varphi^{290}\lP)^{2}
  = \frac{8}{\varphi^{3}}
  \approx 1.889.
\end{equation}
\end{proposition}

\begin{proof}
\begin{align*}
  \Lambda \cdot R^{2}
  &= \frac{2}{\varphi^{583}} \cdot \lP^{-2}
     \cdot 4\varphi^{580} \cdot \lP^{2} \\
  &= \frac{8\varphi^{580}}{\varphi^{583}} \\
  &= \frac{8}{\varphi^{3}} \\
  &= \frac{8}{2\varphi+1} = 1.889\ldots
\end{align*}
This is a number of order unity, as required by the de~Sitter relation
$\Lambda \sim 3/R^{2}$. The effective dark energy fraction is
$\Omega_{\Lambda} = \Lambda R^{2}/3 = 8/(3\varphi^{3}) \approx 0.63$,
consistent with the Planck measurement
$\Omega_{\Lambda} = 0.685 \pm 0.013$.
\end{proof}

\subsection{The Dirac Large Number Relation}

\begin{proposition}\label{prop:dirac}
The Dirac relation $R/\lP \sim \alpha_{G}^{-3/2}$ holds approximately:
\begin{equation}\label{eq:dirac-check}
  \frac{R}{\lP} = 2\varphi^{290}
  \qquad\text{vs.}\qquad
  \alpha_{G}^{-3/2} \sim \varphi^{183 \times 3/2} = \varphi^{274.5}.
\end{equation}
\end{proposition}

\begin{proof}
$R/\lP \sim \varphi^{290}$ and
$\alpha_{G}^{-3/2} \sim \varphi^{274.5}$. The ratio is
$\varphi^{290-274.5} = \varphi^{15.5} \approx 1132$, so the Dirac
relation holds to within a factor of ${\sim}\,10^{3}$---the expected
precision for a scaling argument that ignores corrections of order
$d^{2} = 9$.

In the framework, the exact relation is
$290 = 183 \times (E/d)/(d(|A_{5}|+1)) \times \ldots$, a more complex
identity involving multiple dodecahedral constants. The Dirac relation
is a coarse approximation to this exact lattice identity.
\end{proof}

\subsection{The Hubble Constant}

\begin{corollary}[Hubble Constant]\label{cor:hubble}
The Hubble constant is
\begin{equation}\label{eq:hubble}
  H_{0} = \frac{c}{R}
  = \frac{c}{2\varphi^{290}\lP}
  = \frac{c}{1.306 \times 10^{26}\;\mathrm{m}}
  = 2.298 \times 10^{-18}\;\mathrm{s}^{-1}
  = 70.9\;\mathrm{km/s/Mpc}.
\end{equation}
\end{corollary}

\begin{proof}
$c/R = (2.998 \times 10^{8})/(1.306 \times 10^{26}) =
2.296 \times 10^{-18}\;\mathrm{s}^{-1}$. Converting:
$1\;\mathrm{Mpc} = 3.086 \times 10^{22}\;\mathrm{m}$, so
$H_{0} = 2.296 \times 10^{-18} \times 3.086 \times 10^{22}/10^{3}
= 70.9\;\mathrm{km/s/Mpc}$.

The Planck 2018 measurement gives $H_{0} = 67.4 \pm 0.5$~km/s/Mpc~\cite{planck2018}.
The SH0ES measurement gives $H_{0} = 73.0 \pm 1.0$~km/s/Mpc~\cite{riess2022}.
The framework value $70.9$~km/s/Mpc falls between the two---within
the ``Hubble tension'' window---consistent with $R = c \times t_{0}$
being the comoving radius (not the luminosity distance).
\end{proof}

% ══════════════════════════════════════════════════════════════
\section{The Physical Picture}\label{sec:physical}
% ══════════════════════════════════════════════════════════════

\subsection{The Gear Lattice}

In the axiom framework, space is a fixed lattice of Planck-scale
dodecahedral gears. Each gear has:
\begin{itemize}[nosep]
  \item $d = 3$ neighbors (vertex valency of the dodecahedron),
  \item gear ratio $\varphi$ (the axiom: each gear turns $\varphi$ times
    for one turn of its neighbor),
  \item angular extent $120^{\circ} = |2I|$ degrees between edges,
  \item cone angle $\pi/4$ (koppa$\,\times\,\pi$, where koppa $= 1/4$
    from completing the axiom's square).
\end{itemize}

Adjacent gears rotate in opposite directions (normal configuration).
This cancellation-at-distance is why space appears empty: the net
rotation is zero to leading order.

\subsection{Why $\varphi^{291}$ and Not Infinity?}

The lattice has a natural cutoff at $\varphi^{291}$ Planck lengths.
Beyond this scale, the golden carry rule $\varphi^{2} = \varphi + 1$
forces all higher modes to reduce back into the lattice:
\begin{equation}\label{eq:carry}
  \varphi^{n} = F_{n}\varphi + F_{n-1} \quad (n \geq 291),
\end{equation}
where $F_{n}$ is the $n$-th Fibonacci number. At $n = 291$, the
Fibonacci number $F_{291}$ exceeds the 120-cell capacity of $VE = 600$
effective lattice sites. The universe ``fills up'' at 291 golden breaths.

The ceiling $291 = VE/\chi - d^{2}$ is thus the capacity of the
dodecahedral lattice: 300 vertex--edge sites per Euler unit, minus
$d^{2} = 9$ gravitational degrees of freedom. No golden breath beyond
the 291st can propagate without folding back.

\subsection{The ``$+1$'' as Existence}

The axiom $x^{2} = x + 1$ contains an irreducible ``$+1$'': the
constant term that prevents $\varphi$ from being zero. In the
cosmological context, this ``$+1$'' has three manifestations:
\begin{enumerate}[nosep]
  \item \textbf{The Planck length $\lP$}: the irreducible quantum of
    spatial extent---the ``$+1$'' of the axiom realized as a length.
  \item \textbf{The cosmological constant $\Lambda$}: the ``$+1$'' at
    the round-trip turning point, giving exponent $583 = 2 \times 291 + 1$
    instead of $582 = 2 \times 291$.
  \item \textbf{The universe itself}: the ``$+1$'' is why there is
    something rather than nothing. The axiom $x^{2} = x + 1$ forbids the
    trivial solution $x = 0$. Existence is the constant term.
\end{enumerate}

% ══════════════════════════════════════════════════════════════
\section{Numerical Summary}\label{sec:summary}
% ══════════════════════════════════════════════════════════════

\begin{table}[htbp]
\centering
\caption{Numerical verification of the universe radius.}
\label{tab:verification}
\begin{tabular}{@{}lll@{}}
\toprule
Quantity & Value & Source \\
\midrule
$\varphi$ & $1.6180339887\ldots$ & $(1+\sqrt{5})/2$ \\
$\varphi^{290}$ & $4.0403 \times 10^{60}$ & computed \\
$\lP$ & $1.616255 \times 10^{-35}$ m & CODATA 2018 \\
$R_{\mathrm{theory}} = 2\varphi^{290}\lP$
  & $1.306 \times 10^{26}$ m & Theorem~\ref{thm:main} \\
$R_{\mathrm{theory}}$ & $13.80$ Gly & converted \\
$R_{\mathrm{obs}} = c \cdot t_{0}$ & $13.80 \pm 0.04$ Gly & Planck 2018 \\
\midrule
Relative error & $< 0.04\%$ & within measurement uncertainty \\
\bottomrule
\end{tabular}
\end{table}

\begin{table}[htbp]
\centering
\caption{The exponent 290 from dodecahedral constants.}
\label{tab:290}
\begin{tabular}{@{}lll@{}}
\toprule
Expression & Value & Interpretation \\
\midrule
$VE/\chi$ & $300$ & 120-cell vertices per Euler unit \\
$d^{2}$ & $9$ & Gravitational DOF \\
$VE/\chi - d^{2}$ & $291$ & Structural ceiling (breath count) \\
$N_{\mathrm{ceil}} - 1$ & $290$ & Exponent (fence-post) \\
$V \cdot L_{7}/\chi$ & $290$ & $= 20 \times 29 / 2$ \\
$VE/\chi - E/d$ & $290$ & $= 300 - 10$ \\
\bottomrule
\end{tabular}
\end{table}

% ══════════════════════════════════════════════════════════════
\section{Conclusion}\label{sec:conclusion}
% ══════════════════════════════════════════════════════════════

We have derived the radius of the observable universe from the axiom
$x^{2}=x+1$ and the invariants of the dodecahedron. The formula
\[
  R = 2 \cdot \varphi^{290} \cdot \lP = 13.80\;\text{Gly}
\]
contains:
\begin{itemize}[nosep]
  \item The prefactor $2 = \chi = V - E + F$: the Euler characteristic of
    the closed 2-sphere, giving the topological multiplicity.
  \item The exponent $290 = VE/\chi - d^{2} - 1 = 291 - 1$: the
    structural ceiling minus one, encoding the fence-post count of 291
    golden breaths from Planck to Hubble.
  \item The base $\varphi = (1+\sqrt{5})/2$: the golden ratio, root of
    the axiom.
  \item The unit $\lP$: the Planck length, the irreducible quantum of
    spatial extent.
\end{itemize}

The observable universe spans $\varphi^{291}$ Planck lengths in diameter.
This is not a coincidence or a Dirac large number ``hypothesis'': it is
a consequence of the dodecahedral lattice having $VE/\chi = 300$
effective sites, reduced by $d^{2} = 9$ gravitational degrees of freedom,
giving a capacity of exactly 291 golden breaths.

No free parameters are introduced. The size of the universe is
determined by the topology of the dodecahedron.

% ══════════════════════════════════════════════════════════════
% References
% ══════════════════════════════════════════════════════════════
\begin{thebibliography}{10}

\bibitem{planck2018}
N.~Aghanim \textit{et al.} (Planck Collaboration),
``Planck 2018 results.\ VI.\ Cosmological parameters,''
\textit{Astron.\ Astrophys.}\ \textbf{641}, A6 (2020).

\bibitem{dirac1937}
P.~A.~M.\ Dirac,
``The Cosmological Constants,''
\textit{Nature}\ \textbf{139}, 323 (1937).

\bibitem{codata2018}
E.~Tiesinga, P.~J. Mohr, D.~B. Newell, and B.~N. Taylor,
``CODATA recommended values of the fundamental physical constants: 2018,''
\textit{Rev.\ Mod.\ Phys.}\ \textbf{93}, 025010 (2021).

\bibitem{axiom-theory}
nos3bl33d/DFG,
``The Pythagorean Framework: Icosahedral Structure in Fundamental Constants,''
unpublished manuscript, April 2026.

\bibitem{axiom-lambda}
nos3bl33d/DFG,
``The Cosmological Constant from $x^{2}=x+1$,''
unpublished manuscript, April 2026.

\bibitem{riess2022}
A.~G.~Riess \textit{et al.},
``A Comprehensive Measurement of the Local Value of the Hubble Constant,''
\textit{Astrophys.\ J.\ Lett.}\ \textbf{934}, L7 (2022).

\end{thebibliography}

\end{document}
